% Generated draft: requires mathematical and source review.
\documentclass[11pt]{article}
\usepackage[margin=25mm]{geometry}
\usepackage{fontspec}
\setmainfont{Latin Modern Roman}
\usepackage{amsmath,amssymb,mathtools,unicode-math}
\setmathfont{Latin Modern Math}
\usepackage{enumitem,needspace,xurl,hyperref}
\setlength{\parindent}{0pt}
\setlength{\parskip}{5pt}
\setlength{\emergencystretch}{3em}
\newcommand{\N}{\mathbb{N}}
\newcommand{\Z}{\mathbb{Z}}
\newcommand{\Q}{\mathbb{Q}}
\newcommand{\R}{\mathbb{R}}
\newcommand{\C}{\mathbb{C}}
\newcommand{\F}{\mathbb{F}}
\newcommand{\A}{\mathbb{A}}
\newcommand{\PP}{\mathbb{P}}
\newcommand{\E}{\mathbb{E}}
\newcommand{\eps}{\varepsilon}
\newcommand{\dd}{\,\mathrm{d}}
\begin{document}
\section*{UnsolvedMath ID: 20000427}
\textbf{Source problem code:} catalogue entry
\subsection{Definitions and mathematical statement}
Let \(E/\Q\) be an elliptic curve with complex multiplication by an imaginary quadratic field \(K\), and let \(\chi\) be its associated Hecke character. The source asks for an explicit algebraic correspondence from \(A=E\times E\) to an elliptic modular surface \(\mathcal E\to X_{\mathrm{mod}}\), realizing their common rank-two \(\chi^2\), \(\overline\chi^{,2}\) summand in degree-two cohomology. It associates \(\chi^2\) with a weight-three modular form, and gives its level as the norm of the conductor of \(\chi^2\) times the absolute discriminant of \(K\). A correspondence between these surfaces means a rational codimension-two cycle on \(A\times\mathcal E\), not necessarily a rational map. The target surface, level convention, field and required action on the summand must be specified.

The attached aid reports an explicit source projector. Choose a nonzero trace-zero CM endomorphism \(\delta\), put \(n=\deg\delta\), let \(a=\delta\times\delta\) act on \(A\), and let \(\pi_2\) be its degree-two Chow--K\"unneth projector. With \(\Gamma_a\) the graph, it reports the rational Chow idempotent
\[
p_T=\frac12\left(\pi_2-\frac1n\pi_2\circ{}^t\Gamma_a\circ\pi_2\right).
\]
Its realizations are reported to be exactly the two-dimensional \(\chi^2\), \(\overline\chi^{,2}\) summand. Although \(\delta\) may require a CM field extension, the aid reports descent of the projected correspondence \(p_T\) to \(\Q\); this is not an assertion that \(a\) itself is a \(\Q\)-endomorphism.

Let \(S=\operatorname{Km}(A)\) be the resolved Kummer quotient, and let \(\gamma\) be the resolved degree-two quotient graph regarded as a Chow correspondence. The aid reports
\[
Z=\gamma\circ p_T,\qquad
W=\frac12p_T\circ{}^t\gamma,
\qquad W\circ Z=p_T,\qquad Z\circ W=p_S,
\]
where \(p_S\) is the corresponding transcendental Chow projector on \(S\). These are claimed identities in rational Chow motives, not only cohomology. The source identifies Shioda's level-four example as a K3 Kummer surface associated with \(E\times E\); the aid says this supplies the requested correspondence in that known Kummer case. The level-four surface convention should not be silently identified with every newform-level convention.

For a general elliptic modular surface, the aid only reduces the task to finding a projected cycle
\[
e_f\circ\zeta\circ p_T,
\qquad \zeta\in\mathrm{CH}^2(A\times\mathcal E)_{\Q},
\]
that is nonzero in the intended common-motive realization, where \(e_f\) selects the modular-form summand. Its precise projector and nonvanishing criterion are not supplied. Merely being a nonzero Chow class would not by itself specify a nonzero cohomological action. The original C. Schoen comment relates this to extra cycles after supersingular reduction of \(E\); it is retained as source context, not a construction of the arbitrary modular-surface correspondence.
\subsection{Short English statement}
Construct explicitly the shared CM weight-three correspondence between an elliptic self-product and an elliptic modular surface. The attached aid reports a rational CM projector and inverse Kummer correspondences for the known Kummer case, while leaving the general target to a projected-cycle construction.
\subsection{Sources}
\begin{itemize}
\item Local UnsolvedMath registry, permanent ID \(20000427\), exact public code \texttt{catalogue entry}; canonical location \texttt{aim-algebraic-number-theory-notes.json}, notes[78].
\item AIM workshop ``The Tate conjecture,'' attributed to D. Ramakrishnan; supplied URL \url{https://aimath.org/WWN/tateconjecture/tateconjecture.pdf}. Its verification corrects extracted item 6 to Problem 16, distinguishes elliptic curve \(E\) from elliptic surface \(\mathcal E\), and repairs \(\chi^2\), \(H^2\) and layout. Comments name Shioda's level-four example and C. Schoen's supersingular observation. PDF/TeX were not consulted here; the literature field is empty.
\item Attached research attempt 1, ``An explicit CM projector and Kummer correspondence,'' without complete bibliographic citations or the arbitrary modular-surface projector data.
\end{itemize}
\end{document}
